\section{Datasheet for OenoBench}
\label{app:datasheet}

We follow the datasheet template of \citet{gebru2021datasheets}.

\subsection{Motivation}
\paragraph{For what purpose was the dataset created?}
To evaluate factual knowledge of large language models on a multi-disciplinary,
externally-validated domain (wine), and to provide a benchmark whose
construction includes explicit bias controls (multi-model generation,
multi-agent audit, self-preference scoring, closed-book vs.\ source-grounded
partition).

\paragraph{Who created the dataset?}
Author and affiliation details are withheld for double-blind review.
The dataset was created for the NeurIPS 2026 Evaluations \& Datasets
Track. The lead author holds the WSET Diploma in Wines (the highest
pre-Master-of-Wine qualification of the Wine \& Spirit Education
Trust) and produced the gold-sheet ratings used to calibrate the
audit agents.

\paragraph{Funding.}
Funding source withheld for double-blind review. No external research
grants were received. Total LLM API cost across the entire project
--- generation pilots, audit cycles, and the final 16-config
evaluation --- was \textbf{\$783}, paid out-of-pocket.

\subsection{Composition}
\paragraph{What do the instances represent?}
Each instance is a multiple-choice benchmark question with: question
text, 4~options with marked-correct answer, supporting fact(s) with
source URL and tier label, generator model identifier, generation
strategy, audit-agent verdicts (per-agent pass/warn/fail signals), and
(for the gold-sheet subset) human-reviewer ratings on eight rubrics.

\paragraph{How many instances are there?}
\textbf{3{,}266} questions in \texttt{release\_v1.2}. By strategy:
fact-to-question 1{,}909, distractor mining 405, template 389,
scenario 319, comparative 244. By difficulty (post-relabel):
L1 694, L2 894, L3 678, L4 1{,}001. By domain: wine regions 1{,}108,
grape varieties 766, producers 515, viticulture 502, wine business 250,
winemaking 187. The corpus is built from 38{,}104 atomic facts across
35 scrapers and 580 unique source URLs.

\paragraph{What data does each instance consist of?}
Structured Postgres records mirroring the schema documented in
Section~\ref{sec:dataset} and Appendix~\ref{app:methodology}.

\paragraph{Recommended data splits.}
The dataset ships as a single \texttt{test} split (this is an
evaluation-only benchmark). Two analytic partitions are documented and
recommended for use:
(i) the \emph{closed-book} slice (1{,}601 questions tagged
\texttt{closed\_book\_solvable}) and the \emph{contextual} slice
(1{,}665 remaining questions), partitioned by audit agent B2; and
(ii) per-generator held-out splits for self-preference analysis
(Section~\ref{sec:eval-sps}), defined dynamically by the
\texttt{generator\_family} column.

\paragraph{Errors, noise, redundancies?}
Documented in Section~\ref{sec:audit}. Per-agent fail rates and
$\kappa$ versus the human gold sheet are released alongside the
corpus. The 1{,}601 B2-flagged questions are kept with disclosure
because LLM-judge calibration with humans on closed-book solvability
is unreliable ($\kappa \approx 0.007$).

\paragraph{Self-contained, or relies on external resources?}
Self-contained: question text, supporting facts, and source URLs are
all included. Source URLs may rot over time; we mitigate by archiving
Tier-1 source pages and recording the access timestamp.

\subsection{Collection process}
\paragraph{How was the data acquired?}
By 35 provenance-verified web scrapers (released in
\texttt{src/scrapers/}). Sources span government registries (INAO,
TTB, OIV), inter-governmental bodies, university research groups
(notably UC Davis, USDA Extension), Wikipedia, Wikidata, peer-reviewed
journals (\emph{OENO One}, \emph{Vitis}, \emph{AJEV}), and curated
open datasets. Scraping was rate-limited and used the disclosed user
agent \texttt{OenoBench-Research/1.0 (academic wine benchmark)}.

\paragraph{Time frame.}
Fact extraction: 2026-03 through 2026-04. Question generation:
2026-04 through 2026-05-03. Evaluation: 2026-05-03.

\paragraph{Ethical review.}
The work is non-IRB-eligible at our institution because (a) no human
subjects were enrolled in research, (b) the gold-sheet reviewer is
also the lead author, and (c) all source data is public-record
information about commercial wine entities. Human review of
LLM-generated questions is treated as data-quality assurance, not as
human-subjects research.

\subsection{Preprocessing / cleaning / labelling}
Atomic-fact extraction with entity tagging; per-source paraphrase to
avoid verbatim copy; near-duplicate suppression at cosine~$\geq 0.92$;
nine-agent automated audit; human gold-sheet rating on a stratified
50-question subset per release.

\subsection{Uses}
\paragraph{Recommended uses.} Evaluation of LLM factual knowledge in
the wine domain; ablation studies of bias-mitigation techniques in
benchmark construction; study of self-preference effects in
LLM-as-judge pipelines; calibration of automated multi-agent QA
frameworks against expert human review.

\paragraph{Discouraged uses.} Direct deployment as a wine-recommendation
or purchasing system without separate alignment, age-gating, and
jurisdictional advertising review; fine-tuning on the benchmark itself
for ranking purposes (would invalidate the benchmark for the trained
model). The scrape-to-audit pipeline should not be retargeted to
high-stakes domains (medicine, law) without comparable expert human
review of the retargeted source taxonomy.

\subsection{Distribution}
\paragraph{Will the dataset be distributed?}
Yes, under \textbf{CC-BY-SA-4.0} on HuggingFace; the anonymous review
URL and the camera-ready URL are recorded in
Appendix~\ref{app:artifacts}. A Croissant manifest is included
\citep{koehn2024croissant}. The construction code is hosted in a
double-blind-anonymous mirror also documented in
Appendix~\ref{app:artifacts}.

\paragraph{Citation.}
\begin{verbatim}
@inproceedings{anonymous2026oenobench,
  title  = {OenoBench: A Wine-Domain Benchmark for Knowledge-Grounded
            Evaluation of Large Language Models},
  author = {Anonymous},
  booktitle = {Advances in Neural Information Processing Systems
               (NeurIPS), Evaluations and Datasets Track},
  year   = {2026},
  note   = {Under double-blind review}
}
\end{verbatim}

\subsection{Maintenance}
\paragraph{Maintainer.}
Maintainer details withheld for double-blind review. After acceptance
this section will list the lead author's name, institutional
affiliation, and contact email; during review, all queries are
routed through the OpenReview submission system.

\paragraph{Update cadence.}
Wine regulation changes on a multi-year cadence; we plan annual
re-extraction. Scraper code is released so the community can
re-extract independently. Corpus versioning follows
\texttt{release\_vX.Y} (major release / minor patch) with all prior
versions retained on HuggingFace.

\paragraph{Errata.}
A public errata log will be maintained alongside the HuggingFace
release; corrections to questions, facts, audit verdicts, or
difficulty labels will be versioned and dated. Reports of factual
errors are accepted via GitHub issues and will be evaluated against
the source URL.
